\subsection{Inverse limits of measures}

Let X, Y be two compact spaces, \(p : X \to Y\) a continuous mapping; then \(f \mapsto f \circ p\) is a continuous linear mapping of \(\mathcal{C}(Y; \mathbf{C})\) into \(\mathcal{C}(X; \mathbf{C})\) , since \(\|f \circ p\| \leqslant \|f\|\) for every function \(f \in \mathcal{C}(Y; \mathbf{C})\) ; we denote this mapping by \(p'\) . Its transpose \(^{t}p' : M(X; C) \to M(Y; C)\) is therefore such that, for every measure \(\mu\) on X, \(^{t}p'(\mu)\) is the measure on Y such that

\[
\langle^ {t} p ^ {\prime} (\mu), f \rangle = \langle \mu , f \circ p \rangle
\]

for every function \(f \in \mathcal{C}(\mathrm{Y};\mathbf{C})\) . Note that for every \(x \in \mathrm{X}\) , \(^t p'(\varepsilon_x) = \varepsilon_{p(x)}\) ; for this reason, we shall denote by \(p_*(\mu)\) the measure \(^t p'(\mu)\) , which thus extends \(p\) when \(\mathrm{X}\) (resp. \(\mathrm{Y}\) ) is canonically embedded in \(\mathcal{M}(\mathrm{X};\mathbf{C})\) (resp. \(\mathcal{M}(\mathrm{Y};\mathbf{C})\) ) (§1, No.~9, Prop. 13); for every measure \(\mu\) on \(\mathrm{X}\) , \(p_*(\mu)\) is a special case of the general concept of image of a measure, which we shall study in Ch. V, §6. Since, as we saw above, \(\| p' \| \leqslant 1\) , we have also \(\| ^t p' \| \leqslant 1\) and so

\[
\| p _ {*} (\mu) \| \leqslant \| \mu \|\tag{15}
\]

for every measure \(\mu \in \mathcal{M}(\mathrm{X};\mathbf{C})\)

Let us now consider a directed pre-ordered set I, and an inverse system (or `projective system') \((\mathrm{X}_{\alpha}, p_{\alpha\beta})\) of compact spaces \(X_{\alpha}\) (GT, I, §4, No.~4) having I as set of indices; the inverse limit space \(X = \varprojlim X_{\alpha}\) is known to be compact (GT, I, §9, No.~6, Prop. 8); we shall denote by \(p_{\alpha}\) the canonical mapping of X into \(X_{\alpha}\) .

It is clear that \((\mathcal{M}(\mathrm{X}_{\alpha};\mathbf{C}),(p_{\alpha\beta})_{*})\) is an inverse system of vector spaces, and that \(((p_{\alpha})_{*})\) is an inverse system of linear mappings, which justifies the following definition:

DEFINITION 2. --- A family \((\mu_{\alpha})_{\alpha \in I}\) , where, for every \(\alpha \in I\) , \(\mu_{\alpha}\) is a measure on \(X_{\alpha}\) , is said to be an inverse system of measures if, whenever \(\alpha \leqslant \beta\) , \(\mu_{\alpha} = (p_{\alpha \beta})_{*}(\mu_{\beta})\) . A measure \(\mu\) on \(X = \varprojlim X_{\alpha}\) is said to be an inverse limit of the inverse system \((\mu_{\alpha})\) if \(\mu_{\alpha} = (p_{\alpha})_{*}(\mu)\) for every \(\alpha \in I\) .

PROPOSITION 8. --- (i) If an inverse system \((\mu_{\alpha})\) of measures on the \(X_{\alpha}\) has an inverse limit, then that limit is unique.

\begin{enumerate}
\def\labelenumi{(\roman{enumi})}
\setcounter{enumi}{1}
\item
  If an inverse system \((\mu_{\alpha})\) has an inverse limit, then the family of norms \((\| \mu_{\alpha}\|)\) is bounded.
\item
  If the \(p_{\alpha\beta}\) are surjective and the family ( \(\|\mu_{\alpha}\|\) ) is bounded, then the inverse system of measures ( \(\mu_{\alpha}\) ) has an inverse limit.
\item
  If the \(p_{\alpha \beta}\) are surjective, then every inverse system \((\mu_{\alpha})\) of positive measures has an inverse limit \(\mu\) , which is a positive measure on \(X\) , and \(\| \mu \| = \| \mu_{\alpha} \|\) for all \(\alpha\) .
\end{enumerate}

\begin{enumerate}
\def\labelenumi{(\roman{enumi})}
\tightlist
\item
  We first prove the following lemma:
\end{enumerate}

Lemma 3. --- Let F be the set of functions \(f \in \mathcal{C}(\mathrm{X};\mathbf{C})\) having the following property: there exist an \(\alpha \in \mathrm{I}\) and a function \(f_{\alpha} \in \mathcal{C}(\mathrm{X}_{\alpha};\mathbf{C})\) such that \(f = f_{\alpha} \circ p_{\alpha}\) . Then F is a dense linear subspace of \(\mathcal{C}(\mathrm{X};\mathbf{C})\) .

We note first of all that if \(g = g_{\beta} \circ p_{\beta}\) and \(h = h_{\gamma} \circ p_{\gamma}\) , where \(g_{\beta} \in \mathcal{C}(\mathrm{X}_{\beta}; \mathbf{C})\) and \(h_{\gamma} \in \mathcal{C}(\mathrm{X}_{\gamma}; \mathbf{C})\) , then there exists an \(\alpha \in I\) such that \(\alpha \geqslant \beta\) and \(\alpha \geqslant \gamma\) , and therefore \(p_{\beta} = p_{\beta\alpha} \circ p_{\alpha}\) , \(p_{\gamma} = p_{\gamma\alpha} \circ p_{\alpha}\) , which shows that

\[
g + h = \left(g _ {\beta} \circ p _ {\beta \alpha} + h _ {\gamma} \circ p _ {\gamma \alpha}\right) \circ p _ {\alpha}
\]

belongs to F; one sees similarly that \(gh \in F\) ; F is thus a C-subalgebra of \(\mathcal{C}(\mathrm{X};\mathbf{C})\) , which contains the constants and is such that the relation \(f \in F\) implies \(\overline{f} \in F\) . Finally, if \(x \neq y\) are two points of X, there exists an \(\alpha \in I\) such that \(p_{\alpha}(x) \neq p_{\alpha}(y)\) , therefore there is a function \(f_{\alpha} \in \mathcal{C}(\mathrm{X}_{\alpha};\mathbf{C})\) such that \(f_{\alpha}(p_{\alpha}(x)) \neq f_{\alpha}(p_{\alpha}(y))\) . The conclusion therefore follows from the Stone--Weierstrass theorem (GT, X, §4, No.~4, Prop. 7).

The lemma having been established, let \(\mu, \mu'\) be two measures on X such that \((p_{\alpha})_{*}(\mu) = (p_{\alpha})_{*}(\mu')\) for all \(\alpha \in I\) ; this means that, for every \(\alpha \in I\) and every function \(f_{\alpha} \in \mathcal{C}(X_{\alpha}; \mathbf{C})\) , we have

\[
\left\langle f _ {\alpha}, \left(p _ {\alpha}\right) _ {*} (\mu) \right\rangle = \left\langle f _ {\alpha}, \left(p _ {\alpha}\right) _ {*} \left(\mu^ {\prime}\right) \right\rangle ,
\]

that is, \(\langle f_{\alpha} \circ p_{\alpha}, \mu \rangle = \langle f_{\alpha} \circ p_{\alpha}, \mu' \rangle\) ; in other words, \(\mu\) and \(\mu'\) coincide on the subspace F of \(\mathcal{C}(\mathrm{X}; \mathbf{C})\) , which is dense by Lemma 3, therefore \(\mu = \mu'\) , which proves (i).

\begin{enumerate}
\def\labelenumi{(\roman{enumi})}
\setcounter{enumi}{1}
\tightlist
\item
  The relation (15) applied to \(p_{\alpha}\) shows that if \(\mu\) is the inverse limit of the inverse system \((\mu_{\alpha})\) , necessarily
\end{enumerate}

\[
\| \mu \| \geqslant \| \mu_ {\alpha} \|\tag{16}
\]

for all \(\alpha \in I\) .

\begin{enumerate}
\def\labelenumi{(\roman{enumi})}
\setcounter{enumi}{2}
\tightlist
\item
  Suppose the \(p_{\alpha \beta}\) are surjective; one knows that the same is then true of the \(p_{\alpha}\) (GT, I, §9, No.~6, Prop. 8). Consider an inverse system of measures \((\mu_{\alpha})\) and let us first show that there exists a linear form \(\lambda\) on F (in the notations of Lemma 3) such that, for every \(\alpha \in \mathbf{I}\) and every \(f_{\alpha} \in \mathcal{C}(\mathrm{X}_{\alpha}; \mathbf{C})\) , \(\lambda(f_{\alpha} \circ p_{\alpha}) = \mu_{\alpha}(f_{\alpha})\) . To that end, let \(\beta, \gamma\) be two indices in I, and \(f_{\beta} \in \mathcal{C}(\mathrm{X}_{\beta}; \mathbf{C})\) , \(f_{\gamma} \in \mathcal{C}(\mathrm{X}_{\gamma}; \mathbf{C})\) two functions such that \(f_{\beta} \circ p_{\beta} = f_{\gamma} \circ p_{\gamma}\) ; then there exists an index \(\alpha \in \mathbf{I}\) such that \(\alpha \geqslant \beta\) and \(\alpha \geqslant \gamma\) , therefore
\end{enumerate}

\[
p _ {\beta} = p _ {\beta \alpha} \circ p _ {\alpha}, p _ {\gamma} = p _ {\gamma \alpha} \circ p _ {\alpha} \quad \text {and} \quad (f _ {\beta} \circ p _ {\beta \alpha}) \circ p _ {\alpha} = (f _ {\gamma} \circ p _ {\gamma \alpha}) \circ p _ {\alpha};
\]

since \(p_{\alpha}\) is surjective, this implies \(f_{\beta} \circ p_{\beta \alpha} = f_{\gamma} \circ p_{\gamma \alpha}\) , therefore

\[
\mu_ {\alpha} (f _ {\beta} \circ p _ {\beta \alpha}) = \mu_ {\alpha} (f _ {\gamma} \circ p _ {\gamma \alpha});
\]

but by definition the last relation may also be written \(\mu_{\beta}(f_{\beta}) = \mu_{\gamma}(f_{\gamma})\) , whence our assertion.

This being so, suppose that there exists a finite number \(a > 0\) such that \(\| \mu_{\alpha}\| \leqslant a\) for all \(\alpha \in \mathrm{I}\) ; then, for every function \(f_{\alpha} \in \mathcal{C}(\mathrm{X}_{\alpha};\mathbf{C})\) ,

\[
\left| \lambda (f _ {\alpha} \circ p _ {\alpha}) \right| = \left| \mu_ {\alpha} (f _ {\alpha}) \right| \leqslant a \| f _ {\alpha} \| = a \| f _ {\alpha} \circ p _ {\alpha} \|
\]

since \(p_{\alpha}\) is surjective. This shows that the linear form \(\lambda\) is continuous on \(\mathbf{F}\) , and it follows from Lemma 3 that \(\lambda\) may be extended to a measure \(\mu\) on \(\mathbf{X}\) such that \((p_{\alpha})_{*}(\mu) = \mu_{\alpha}\) for all \(\alpha \in \mathbf{I}\) , which proves (iii).

\begin{enumerate}
\def\labelenumi{(\roman{enumi})}
\setcounter{enumi}{3}
\tightlist
\item
  To prove the existence of \(\mu\) it suffices, by (iii), to verify that the family of norms ( \(\| \mu_{\alpha}\|\) ) is bounded. But \(\| \mu_{\alpha}\| = \mu_{\alpha}(1)\) and, when \(\alpha \leqslant \beta\) , the relation \(\mu_{\alpha} = (p_{\alpha \beta})_{*}(\mu_{\beta})\) implies that \(\mu_{\alpha}(1) = \mu_{\beta}(1)\) ; since I is directed, the total masses of all the measures \(\mu_{\alpha}\) are therefore equal, whence our assertion. Moreover, the subspace F obviously satisfies the property (P) of §1, No.~7, Prop. 9, therefore the inverse limit measure \(\mu\) of \((\mu_{\alpha})\) is positive. Finally, the relation \(\mu_{\alpha} = (p_{\alpha})_{*}(\mu)\) shows as above that \(\mu(1) = \mu_{\alpha}(1)\) .
\end{enumerate}

Example. --- Let \((\mathrm{X}_{\lambda})_{\lambda\in\mathrm{L}}\) be a family of compact spaces; set \(X=\prod_{\lambda\in L}X_{\lambda}\) and, for every finite subset J of L, set \(X_{J}=\prod_{\lambda\in J}X_{\lambda}\) ; denote by \(pr_{J}:X\to X_{J}\) and \(pr_{J,K}:X_{K}\to X_{J}\) (for \(J\subset K\) ) the canonical projections. We know that \((\mathrm{X}_{\mathrm{J}},\mathrm{pr}_{\mathrm{JK}})\) is an inverse system of compact spaces, and that the inverse limit of the system of continuous mappings \((\mathrm{pr}_{\mathrm{J}})\) is a homeomorphism of X onto the inverse limit space \(\varprojlim X_{J}\) , permitting one to identify these two spaces (GT, I, §4, No.~4 and S, III, §7, No.~2, Remark 3). Since the projections \(pr_{J,K}\) are surjective, it follows from Prop. 8 that the set \(\mathcal{M}(\mathrm{X};\mathbf{C})\) (resp. \(\mathcal{M}_{+}(\mathrm{X})\) ) may be identified with the set of inverse systems \((\mu_{\mathrm{J}})\) such that the family of norms \((\|\mu_{\mathrm{J}}\|)\) is bounded (resp. such that the \(\mu_{J}\) are all positive, and necessarily of the same total mass).

Let us consider in particular the case where, for each \(\lambda \in L\) , a measure \(\mu_{\lambda}\) is given on \(X_{\lambda}\) and one sets \(\mu_{J} = \bigotimes_{\lambda \in J} \mu_{\lambda}\) . If \(J \subset K\) are two finite subsets of I then, for every function \(f_{J} \in \mathcal{C}(X_{J}; \mathbf{C})\) we have, by virtue of formula (14) of No.~4,

\[
\mu_ {K} \left(f _ {J} \circ \operatorname{pr} _ {J, K}\right) = \mu_ {J} \left(f _ {J}\right) \cdot \prod_ {\lambda \in K - J} \mu_ {\lambda} (1).\tag{17}
\]

For \((\mu_{\mathrm{J}})\) to be an inverse system of measures, it is therefore necessary and sufficient that either \(\mu_{\lambda}=0\) for all \(\lambda\in L\) or \(\mu_{\lambda}(1)=1\) for all \(\lambda\in L\) .

